% section: シグマの計算

  \section{シグマの計算}
    ここで1つ新たな概念を導入する.
    ギリシャ文字の$\sum$をシグマといい,シグマを用いてある実数（複素数でも構わない）の集合$\{a_k\}$に対して,\,$x_1$から$x_n$までの総和を
    \[
      \sum_{k=1}^{n} a_k=a_1+a_2+\cdots+ a_n
    \]
    と表す.
    この時以下の事柄が一般に成り立つ.
    \begin{align*}
      &\sum_{k=1}^{n} i=ni\\
      &\sum_{k=1}^{n} ia_k=i\sum_{k=1}^{n} a_k\\
      &\sum_{k=1}^{n} a_k+b_k=\sum_{k=1}^{n} a_k+\sum_{k=1}^{n} b_k
    \end{align*}
    また,一般に以下が成り立つ.
    \begin{align*}
      &\sum_{k=1}^{n} k=\frac{1}{2}n(n+1)\\
      &\sum_{k=1}^{n} k^2=\frac{1}{6}n(n+1)(2n+1)\\
      &\sum_{k=1}^{n} k^3=\left\{\frac{1}{2}n(n+1)\right\}^2\\
    \end{align*}

